\documentclass[aspectratio=169,10pt]{beamer}
\usetheme{skd}

\title[Pertemuan 15 -- Isu Implementasi]{Isu Implementasi Pengontrol Digital}
\subtitle{Implementation Issues of Digital Controllers}
\author{Ir. Estiyanti Ekawati, M.T., Ph.D. \quad\textbar\quad Faqihza Mukhlish, S.T., M.T., Ph.D.}
\institute{Program Studi Teknik Fisika, Fakultas Teknologi Industri, Institut Teknologi Bandung}
\date{Pertemuan 15 \textbar\ Semester I 2026/2027}

\begin{document}

\begin{frame}[plain,noframenumbering]
	\titlepage
\end{frame}

\begin{frame}{Tujuan pertemuan}
	\framesubtitle{Learning objectives}
	\begin{columns}[T]
		\begin{column}{0.55\textwidth}
			\begin{enumerate}
				\item Menganalisis pengaruh tunda komputasi dan \textit{jitter}.
				\item Menjelaskan efek kuantisasi koefisien dan aritmetika titik tetap.
				\item Menyusun struktur program waktu nyata yang aman.
				\item Menerapkan alur uji dari simulasi ke perangkat keras.
				\item Mempresentasikan hasil tugas besar.
			\end{enumerate}
		\end{column}
		\begin{column}{0.4\textwidth}
			\begin{block}{Sub-CPMK S14}
				\small Mengevaluasi isu implementasi pengontrol digital pada mikroprosesor dan mempresentasikan hasil rancangan.\\[3pt]
				\en{Evaluate implementation issues of digital controllers on microprocessors and present design results.}
			\end{block}
			{\small Bacaan: [FR] 8.4--8.5; [PH] Bab 12--13}
		\end{column}
	\end{columns}
\end{frame}

% ===========================================================================
\section{Pewaktuan\texorpdfstring{\\[4pt]{\Large\en{Timing}}}{ / Timing}}
% ===========================================================================

\begin{frame}{Tunda komputasi}
	\framesubtitle{Computational delay}
	\begin{columns}[c]
		\begin{column}{0.55\textwidth}
			\centering
			\begin{tikzpicture}[xscale=1.6,yscale=0.9]
				\draw[-{Stealth}] (0,0) -- (4.4,0) node[right,font=\scriptsize]{$t$};
				\foreach \k in {0,1,2,3} {
						\draw (\k*1.4,0.1) -- (\k*1.4,-0.1) node[below,font=\scriptsize]{$\k T$};
						\fill[skdblue] (\k*1.4,0.35) circle (1.5pt);
						\fill[skdorange!60] (\k*1.4,0.6) rectangle (\k*1.4+0.45,0.85);
						\draw[-{Stealth},skdgreen,thick] (\k*1.4+0.45,0.6) -- (\k*1.4+0.45,0.15);
					}
				\node[font=\scriptsize,text=skdblue,anchor=east] at (-0.05,0.35) {baca $y$};
				\node[font=\scriptsize,text=skdorange,anchor=east] at (-0.05,0.75) {hitung};
				\node[font=\scriptsize,text=skdgreen,anchor=west] at (0.5,0.25) {kirim $u$};
				\draw[decorate,decoration={brace,amplitude=3pt}] (0,0.95) -- node[above,font=\scriptsize]{$T_c$} (0.45,0.95);
			\end{tikzpicture}
		\end{column}
		\begin{column}{0.41\textwidth}
			\small
			Model sebenarnya: $u$ berubah $T_c$ setelah $y$ dibaca. Pilihan:
			\begin{enumerate}
				\item \textbf{Kirim segera}: $T_c \ll T$, abaikan atau perlakukan sebagai tunda kecil.
				\item \textbf{Kirim di cacah berikutnya}: tunda tetap tepat $z^{-1}$, masukkan ke model saat merancang.
				\item \textbf{Pisah komputasi}: hitung bagian yang hanya bergantung data lama \textit{sebelum} cacah, sisakan sedikit operasi setelah membaca $y$.
			\end{enumerate}
		\end{column}
	\end{columns}
\end{frame}

\begin{frame}{Dampak tunda satu cacah yang tidak dimodelkan}
	\framesubtitle{Effect of an unmodelled one-sample delay}
	\begin{columns}[c]
		\begin{column}{0.56\textwidth}
			\centering
			\begin{tikzpicture}
				\begin{axis}[skd,width=8.4cm,height=5.2cm,xmin=0,xmax=4,ymin=0,ymax=1.5,xlabel={$t$ (s)},ylabel={$y$},legend style={at={(0.98,0.05)},anchor=south east}]
					\addplot[skdgray,dashed,thin,domain=0:4,forget plot] {1};
					\addplot[skdblue,const plot] table[x=t,y=y0] {data/p15_delay.dat};
					\addplot[skdred,const plot] table[x=t,y=y1] {data/p15_delay.dat};
					\legend{sesuai desain, + tunda $z^{-1}$}
				\end{axis}
			\end{tikzpicture}
		\end{column}
		\begin{column}{0.4\textwidth}
			\small
			Kompensator \textit{lead} Pertemuan 9 pada plant posisi:\\[3pt]
			\renewcommand{\arraystretch}{1.2}
			\begin{tabular}{@{}l c c@{}}
				\toprule
				 & desain & + tunda \\
				\midrule
				$M_p$ & 16\% & 38\% \\
				$t_s$ & 2,1 s & 3,8 s \\
				$|z|_{\max}$ & 0,82 & 0,90 \\
				\bottomrule
			\end{tabular}\\[4pt]
			Tunda satu cacah menambah fasa $-\omega T$ (dua kali tunda ZOH). Bila $u$ dikirim di cacah berikutnya, \textbf{rancang ulang} dengan $z^{-1}G(z)$.
		\end{column}
	\end{columns}
\end{frame}

\begin{frame}{\textit{Jitter} dan cara mengukurnya}
	\framesubtitle{Jitter and how to measure it}
	\begin{columns}[T]
		\begin{column}{0.5\textwidth}
			\small
			\term{\textit{Jitter}}{jitter}: variasi waktu antar-cacah atau antar-keluaran. Model diskrit mengasumsikan $T$ konstan.\\[3pt]
			Penyebab di Arduino:
			\begin{itemize}
				\item \texttt{Serial.print} panjang di dalam kalang
				\item komunikasi I$^2$C (DAC) dan \texttt{analogRead} berulang
				\item cabang program dengan durasi berbeda
				\item \texttt{delay()} alih-alih pewaktuan \texttt{millis()}
			\end{itemize}
		\end{column}
		\begin{column}{0.46\textwidth}
			\begin{block}{Praktik baik}
				\small
				\begin{itemize}
					\item Ukur \texttt{DeltaTime} dan waktu eksekusi dengan \texttt{micros()} (Modul 0).
					\item Waktu eksekusi terburuk $< 0{,}5T$.
					\item Cetak data lebih jarang atau dalam format biner ringkas.
					\item Untuk $T$ kecil dan presisi tinggi: \textit{timer interrupt}.
				\end{itemize}
			\end{block}
		\end{column}
	\end{columns}
\end{frame}

% ===========================================================================
\section{Numerik\texorpdfstring{\\[4pt]{\Large\en{Numerics}}}{ / Numerics}}
% ===========================================================================

\begin{frame}{Kuantisasi koefisien}
	\framesubtitle{Coefficient quantisation}
	\begin{columns}[T]
		\begin{column}{0.52\textwidth}
			\small
			Filter orde 2 dengan kutub $0{,}98\,e^{\pm j0{,}05}$:
			\[
				z^2 + a_1z + a_2,\ \ a_1 = -1{,}957551,\ \ a_2 = 0{,}960400
			\]
			Koefisien dibulatkan ke $b$ bit pecahan:\\[3pt]
			\renewcommand{\arraystretch}{1.2}
			\begin{tabular}{@{}c c c@{}}
				\toprule
				Bit pecahan & $|z|$ & sudut (rad) \\
				\midrule
				ideal & 0,98000 & 0,0500 \\
				8 & 0,98027 & 0,0599 \\
				12 & 0,98003 & 0,0508 \\
				16 & 0,98000 & 0,0500 \\
				\bottomrule
			\end{tabular}\\[3pt]
			Dengan 8 bit, frekuensi osilasi bergeser 20\%.
		\end{column}
		\begin{column}{0.44\textwidth}
			\begin{alertblock}{Mengapa peka?}
				\small
				Kutub dekat $z=1$ (karena $T$ kecil dibanding dinamika) membuat koefisien hampir saling meniadakan; galat kecil pada $a_i$ menggeser akar jauh.
			\end{alertblock}
			\small
			\textbf{Mitigasi}
			\begin{itemize}
				\item Realisasi kaskade/paralel orde 1--2, bukan bentuk langsung orde tinggi.
				\item Bentuk ruang keadaan atau operator delta $\delta = (z-1)/T$.
				\item Jangan memilih $T$ jauh lebih kecil dari perlu.
			\end{itemize}
		\end{column}
	\end{columns}
\end{frame}

\begin{frame}{Titik tetap vs titik mengambang}
	\framesubtitle{Fixed point vs floating point}
	\begin{columns}[T]
		\begin{column}{0.5\textwidth}
			\small
			\begin{itemize}
				\item ATmega328P (Arduino UNO) tidak punya unit \textit{floating-point}: \texttt{float} diemulasi perangkat lunak, jauh lebih lambat daripada operasi bilangan bulat. Ukur dengan \texttt{micros()}.
				\item Titik tetap (format Q): bilangan bulat dengan skala $2^{-b}$. Cepat, tetapi rentang terbatas: waspadai \textit{overflow}, gunakan aritmetika jenuh.
				\item \texttt{float} AVR = \texttt{double} 32-bit: sekitar 7 digit desimal.
			\end{itemize}
		\end{column}
		\begin{column}{0.46\textwidth}
			\begin{exampleblock}{Contoh: $a = e^{-T/\tau}$ dalam Q8}
				\scriptsize
				\renewcommand{\arraystretch}{1.2}
				\begin{tabular}{@{}c c c c@{}}
					\toprule
					$T$ (s) & $a$ & $a_{Q8}$ & $\tau$ efektif \\
					\midrule
					0,1 & 0,818731 & 0,820313 & 0,505 s \\
					0,01 & 0,980199 & 0,980469 & 0,507 s \\
					0,001 & 0,998002 & 0,996094 & 0,255 s \\
					\bottomrule
				\end{tabular}\\[3pt]
				\small Dengan $T$ sangat kecil, 8 bit pecahan membuat model kehilangan separuh konstanta waktunya.
			\end{exampleblock}
		\end{column}
	\end{columns}
\end{frame}

% ===========================================================================
\section{Praktik Implementasi\texorpdfstring{\\[4pt]{\Large\en{Implementation Practice}}}{ / Implementation Practice}}
% ===========================================================================

\begin{frame}[fragile]{Struktur program waktu nyata}
	\framesubtitle{Real-time program structure}
	\begin{columns}[T]
		\begin{column}{0.54\textwidth}
			\begin{block}{}
				\scriptsize
\begin{verbatim}
void loop() {
  baca_perintah_serial();          // non-blocking
  unsigned long now = micros();
  if (now - last >= T_US) {
    last += T_US;                  // tanpa drift
    float y = baca_sensor();       // 1. ukur
    float u = kontroler(r, y);     // 2. hitung
    u = batasi(u);                 // 3. saturasi
    run_motor(kompensasi(u));      // 4. kirim
    update_keadaan(u);             // 5. pengamat, AW
    if (++n % 5 == 0) log_data();  // 6. catat jarang
  }
}
\end{verbatim}
			\end{block}
		\end{column}
		\begin{column}{0.42\textwidth}
			\small
			\begin{itemize}
				\item Urutan: ukur $\rightarrow$ hitung $\rightarrow$ kirim, lalu pekerjaan lain.
				\item \texttt{last += T\_US} menjaga rata-rata periode tepat $T$.
				\item Keselamatan: batas $u$, deteksi sensor di luar rentang, perintah berhenti darurat.
				\item Inisialisasi semua keadaan (integrator, pengamat) di \texttt{setup()} dan saat berganti mode.
			\end{itemize}
		\end{column}
	\end{columns}
\end{frame}

\begin{frame}{Nonlinieritas yang harus ditangani}
	\framesubtitle{Nonlinearities to handle}
	\begin{columns}[T]
		\begin{column}{0.32\textwidth}
			\begin{block}{Saturasi}
				\small
				Batasi $u$ dan $\Delta u$; \textit{anti-windup} pada integrator dan pengamat (pakai $u$ setelah saturasi).
			\end{block}
		\end{column}
		\begin{column}{0.32\textwidth}
			\begin{block}{\textit{Deadzone} dan gesekan}
				\small
				Kompensasi invers \textit{deadzone} (Modul 1) agar plant mendekati linier di sekitar nol.
			\end{block}
		\end{column}
		\begin{column}{0.32\textwidth}
			\begin{block}{Sensor}
				\small
				Filter anti-alias, kalibrasi, satuan konsisten; perhatikan potensiometer yang ``melompat'' $100\% \leftrightarrow 0\%$.
			\end{block}
		\end{column}
	\end{columns}
	\vspace{0.3cm}
	\kunci{Rancangan linier mengasumsikan semua ini sudah ditangani. Sebagian besar kegagalan implementasi berasal dari sini, bukan dari teori kontrolnya.}
\end{frame}

\begin{frame}{Alur uji: dari simulasi ke perangkat keras}
	\framesubtitle{Test workflow: from simulation to hardware}
	\centering
	\begin{tikzpicture}[node distance=0.35cm and 0.45cm]
		\node[blk,text width=2.1cm] (a) {1. Simulasi model\\(MATLAB/Python)};
		\node[blk,text width=2.1cm,right=of a] (b) {2. Simulasi kode\\kontroler yang sama};
		\node[blk,text width=2.1cm,right=of b] (c) {3. Kit tanpa daya motor\\(cek \textit{Serial Monitor})};
		\node[blk,text width=2.1cm,right=of c] (d) {4. Kit, penguatan kecil, \textit{setpoint} kecil};
		\node[blkd,text width=2.1cm,right=of d] (e) {5. Uji penuh + gangguan, catat data};
		\draw[sig] (a) -- (b); \draw[sig] (b) -- (c); \draw[sig] (c) -- (d); \draw[sig] (d) -- (e);
	\end{tikzpicture}

	\vspace{0.4cm}
	\begin{columns}[T]
		\begin{column}{0.48\textwidth}
			\small
			\begin{itemize}
				\item Langkah 2 menangkap kesalahan satuan, tanda, dan urutan operasi sebelum menyentuh motor.
				\item Langkah 3 adalah SOP Modul 0: lepas transistor atau jangan sambungkan 220\,V.
			\end{itemize}
		\end{column}
		\begin{column}{0.48\textwidth}
			\small
			\begin{itemize}
				\item Bandingkan data kit dengan simulasi pada \textbf{grafik yang sama}.
				\item Bila berbeda jauh: periksa model (identifikasi ulang), tunda, saturasi, \textit{deadzone}.
			\end{itemize}
		\end{column}
	\end{columns}
\end{frame}

% ===========================================================================
\section{Presentasi Tugas Besar dan Rangkuman\texorpdfstring{\\[4pt]{\Large\en{Project Presentations and Course Recap}}}{ / Project Presentations and Course Recap}}
% ===========================================================================

\begin{frame}{Struktur presentasi tugas besar}
	\framesubtitle{Project presentation structure}
	\begin{columns}[T]
		\begin{column}{0.5\textwidth}
			\small
			\textbf{10 menit presentasi + 5 menit tanya jawab}
			\begin{enumerate}
				\item Tujuan dan spesifikasi kinerja (1 slide)
				\item Model dan identifikasi sistem
				\item Metode kontrol: teori singkat dan diagram blok
				\item Perancangan dan simulasi
				\item Implementasi pada kit: hasil, grafik, metrik
				\item Perbandingan dengan PID (Modul 2/3)
				\item Kesimpulan dan pelajaran
			\end{enumerate}
		\end{column}
		\begin{column}{0.46\textwidth}
			\begin{block}{Kriteria penilaian}
				\small
				\renewcommand{\arraystretch}{1.15}
				\begin{tabular}{@{}l r@{}}
					Kelengkapan dokumen & 30\% \\
					Pemahaman sistem \& metode & 30\% \\
					Hasil pengontrolan & 10\% \\
					Presentasi & 30\% \\
					Bonus (kreativitas, demo, uji gangguan) & +15\% \\
				\end{tabular}
			\end{block}
			{\small Dokumen (program, data, laporan) dikumpulkan sebelum UAS.}
		\end{column}
	\end{columns}
\end{frame}

\begin{frame}{Peta metode semester ini}
	\framesubtitle{Map of this semester's methods}
	\centering
	\small
	\renewcommand{\arraystretch}{1.35}
	\begin{tabular}{@{}l l l l@{}}
		\toprule
		 & \textbf{Model} & \textbf{Metode} & \textbf{Kapan dipakai} \\
		\midrule
		Heuristik & FOPDT dari data & PID + ZN/Cohen--Coon ($\theta + T/2$) & cepat, model kasar \\
		Emulasi & $G(s)$ & Tustin/MPZ dari $C(s)$ & $T$ kecil, intuisi kontinu \\
		Langsung $z$ & $G(z)$ (ZOH) & \textit{root locus}, bidang $w$ & $T$ sedang, SISO \\
		Langsung $z$ & $G(z)$ & Ragazzini, \textit{deadbeat}, Dahlin & proses lambat, $T$ besar \\
		Ruang keadaan & $\Phi, \Gamma, C$ & penempatan kutub + pengamat & keadaan internal penting \\
		Optimal & $\Phi, \Gamma$ + derau & LQR, Kalman, LQG, MPC & kompromi, banyak masukan, batasan \\
		\bottomrule
	\end{tabular}

	\vspace{0.3cm}
	\kunci{Semua metode bermuara pada persamaan beda yang dijalankan tiap $T$ detik di mikroprosesor.}
\end{frame}

\begin{frame}{Ringkasan dan persiapan UAS}
	\framesubtitle{Summary and final exam preparation}
	\begin{columns}[T]
		\begin{column}{0.48\textwidth}
			\begin{block}{Ringkasan}
				\small
				\begin{itemize}
					\item Tunda komputasi: kirim segera atau modelkan sebagai $z^{-1}$.
					\item Jaga $T$ tetap; ukur \textit{jitter}.
					\item Koefisien dekat $z=1$ peka kuantisasi; gunakan realisasi orde rendah.
					\item Uji bertahap; tangani saturasi, \textit{deadzone}, dan sensor.
				\end{itemize}
			\end{block}
		\end{column}
		\begin{column}{0.48\textwidth}
			\begin{exampleblock}{Cakupan UAS}
				\small
				Pertemuan 1--15, dengan penekanan pada Pertemuan 9--15: \textit{root locus} bidang $z$, respons frekuensi dan \textit{deadbeat}, ruang keadaan, penempatan kutub, pengamat, LQR/Kalman, dan isu implementasi.
			\end{exampleblock}
			{\small Latihan: ulangi contoh-contoh slide dengan $T = 0{,}05$\,s dan bandingkan hasilnya.}
		\end{column}
	\end{columns}
\end{frame}

\end{document}
